\chapter*{Abstract}
\label{cha:abstract}

Precision agriculture relies heavily on the accurate, large-scale mapping of
land use and crop distribution derived from Earth Observation satellites,
particularly the high-resolution, high-frequency multispectral imagery provided by the Sentinel-2 mission.
By monitoring how specific agricultural parcels reflect light over time,
researchers construct Satellite Image Time Series (SITS)
to track the unique phenological growth stages of varying vegetation.
Despite these advancements, employing deep learning for SITS on a national or global
scale presents severe technical and agronomic challenges.
Satellite imagery is inherently irregular due to persistent cloud cover and atmospheric interference,
and processing high-dimensional spatial-temporal grids creates critical memory and
input/output (I/O) bottlenecks in High-Performance Computing (HPC) environments.
Furthermore, many current state-of-the-art sequence models suffer from ``temporal myopia''.
They treat each growing season as an entirely isolated event, attempting
to classify crops based on just one year of weather and growth patterns.
In doing so, they completely ignore a fundamental rule of farming: crop rotation.
Farmers intentionally alternate specific crops across consecutive
years to maintain soil health and prevent nutrient depletion.
If a satellite observation is heavily obscured by clouds or spectrally ambiguous in a given year,
a model suffering from temporal myopia has no historical context to help it classify the vegetation.
Conversely, a model capable of tracking multi-year rotations can analyze what the farmer planted
in the previous and subsequent years to logically deduce the missing crop.

This thesis proposes a highly scalable machine learning framework designed to solve these
hardware limitations while simultaneously capturing both
intra-year phenology and inter-year crop rotation dynamics.
To overcome the severe HPC bottlenecks associated with large-scale spatial processing,
the framework transitions the data pipeline to a highly optimized \textbf{1D pixel-based extraction method}.
By reducing the data representation to the deep temporal evolution of individual pixels,
the architecture bypasses memory limits and enables the rapid processing of
a massive six-year continuous SITS dataset over Austria (INVEKOS).
To support rigorous architectural evaluation, the project utilizes an automated
Machine Learning Operations (MLOps) pipeline integrating PyTorch Lightning, Ray Tune,
and Optuna to guarantee fair hyperparameter tuning.
Additionally, to handle irregular temporal gaps specifically within the Transformer models,
the framework introduces a custom \textbf{AgroDays Positional Encoding} module, which
scales sinusoidal encodings using the exact agricultural day of the satellite acquisition
to explicitly inject temporal distance into the spectral embeddings.

The research systematically establishes robust single-year baselines to serve as
standard benchmarks, including a \textbf{1-Dimensional Convolutional Neural Network (1D-CNN)}
and standalone \textbf{Transformer Encoders} evaluated in both generative and discriminative variants.
Crucially, the research also introduces \textbf{linear structural controls};
in these tests, the deep non-linear feature extractors are replaced by a simplistic linear
backbone that feeds directly into multi-year sequence heads.
This deliberate structural control definitively proves the mathematical necessity of
utilizing high-capacity, non-linear neural backbones before multi-year sequence smoothing
can be applied to raw satellite data. 

To bridge the gap between phenological tracking and historical agricultural dependencies,
two distinct multi-year hybrid architectures are confronted.
The established \textbf{Transformer-HMM cascade} acts as the Bayesian reference,
where the Transformer estimates yearly class probabilities and a Hidden Markov
Model (HMM) smooths these predictions over time using explicit, fixed transition matrices.
Against this, the thesis proposes a novel \textbf{Hierarchical Transformer-LSTM},
a purely neural end-to-end architecture.
It processes intra-year sequences in parallel using self-attention, and feeds the chronological
yearly predictions into an inter-year bidirectional LSTM.
This forces the network to learn multi-year crop rotation rules from past and future planting
contexts without the use of predefined agronomic matrices.

Experimental results across 47 distinct crop classes demonstrate the severe limitations of
single-year models on highly imbalanced datasets.
The 1D-CNN and standalone Discriminative Transformer achieved limited Macro F1-scores
of 30.17\% and 33.18\%, respectively, while the linear structural controls suffered
catastrophic performance degradation.
Ultimately, both multi-year architectures vastly outperformed the single-year baselines,
demonstrating that historical crop rotation context is highly effective at
rescuing minority classes that traditional models ignore.
The novel Transformer-LSTM successfully learned complex agricultural cycles end-to-end,
achieving classification performance competitive with the state-of-the-art HMM cascade.
By delivering this optimized framework, this thesis provides a highly scalable software
solution for operational crop mapping, laying the groundwork for future research which might
integrate State Space Models, Graph Neural Networks, and multimodal SAR data fusion.